\documentclass[12pt, titlepage]{article}

\usepackage{amsmath, mathtools}
\usepackage[round]{natbib}
\usepackage{amsfonts}
\usepackage{amssymb}
\usepackage{graphicx}
\usepackage{hyperref}
\usepackage{longtable}
\usepackage{tabularx}
\usepackage{float}
\usepackage{booktabs}
\usepackage{caption}
\usepackage{fullpage}
\usepackage{array}

\hypersetup{
bookmarks=true,
colorlinks=true,
linkcolor=red,
citecolor=blue,
filecolor=magenta,
urlcolor=cyan
}

\input{../../Comments.text}
\commentsfalse
\input{../../Common.text}

\begin{document}

\title{Module Interface Specification for \progname{}}
\author{\authname}
\date{\today}

\maketitle

\pagenumbering{roman}

\section{Revision History}

\begin{tabularx}{\textwidth}{p{3cm}p{2cm}X}
\toprule {\bf Date} & {\bf Version} & {\bf Notes}\\
\midrule
Nov 13, 2026 & 0.1 & Tyler Yue: Drafted Rev~-1 and ensured MIS syntax for all leaf modules
aligned with the Module Guide.\\
Nov 13, 2026 & 0.2 & Tyler Yue: Edited the Rev~-1 to ensure it is aligned with the rubric and the Module Guide.\\
Nov 16, 2026 & 0.3 & Vedant Patel: Added reflection appendix with individual reflection.\\
\bottomrule
\end{tabularx}

~\newpage

\section{Symbols, Abbreviations and Acronyms}

See also the SRS (\texttt{docs/SRS-Volere/SRS.tex}) and Module Guide
(\texttt{docs/Design/SoftArchitecture/MG.tex}).

\begin{center}
\begin{tabular}{l p{10cm}}
\toprule
\textbf{Symbol} & \textbf{Description}\\
\midrule
MG & Module Guide\\
MIS & Module Interface Specification\\
AO & Abstract Object\\
FR & Functional Requirement\\
JML & Java Modeling Language\\
LLM & Large Language Model\\
$\mathbb{B}$ & Boolean\\
$\mathbb{Z}$ & Integer\\
$\mathbb{N}$ & Natural number\\
seq of $X$ & Sequence of values of type $X$\\
\bottomrule
\end{tabular}
\end{center}

\newpage

\tableofcontents

\newpage

\pagenumbering{arabic}

\section{Introduction}

This document gives the Module Interface Specifications for \progname{}, a
VS~Code extension that captures engineering intent, requires human approval of
proposed artifacts, stores a versioned intent model with rationale and
traceability, and verifies AI-generated changes against that model.

Complementary documents include the Software Requirements Specification, Hazard
Analysis, Module Guide, and VnV Plan. The repository is
\url{https://github.com/RishabShine/IDE3.0}.

This submission is \textbf{Design Documentation Revision~-1}. Per the course
requirements, this revision emphasizes \textbf{MIS syntax}: access program
names, arguments, return types, state variables, environment variables, exported
constants, and exceptions. Full formal access-routine semantics, and design
reflection, are planned for Revision~0.

\subsection{Interface Design Qualities}

Access programs are written to satisfy the Hoffman--Strooper interface
qualities~\citep{HoffmanAndStrooper1995}, applied as follows:

\begin{itemize}
  \item \textbf{Consistency.} Naming and parameter order follow the same pattern
    across modules (identifier of the request or artifact first; domain types
    rather than raw character sequences for structured values).
  \item \textbf{Essentiality.} Each program corresponds to a service stated in
    the Module Guide; constructor-style or framework lifecycle names are not
    exported.
  \item \textbf{Generality.} Programs speak in terms of Request, Proposal,
    Decision, ContractSet, and related derived types so a future host or
    assistant binding does not change the MIS signatures.
  \item \textbf{Minimality.} Leaf modules export only the operations needed by
    their users in the MG uses hierarchy; convenience wrappers that duplicate
    another module's secret are omitted.
  \item \textbf{High cohesion / low coupling.} State owned by a module is
    accessible only through that module's programs; uses relations match
    the Module Guide uses hierarchy and introduce no cycles.
  \item \textbf{Opacity.} Encodings (file layout, host widget APIs, assistant
    message formats, checker invocation) remain secrets of the owning module;
    clients never construct those encodings.
\end{itemize}

Abstraction is kept above implementation syntax: the MIS does not prescribe
serialization formats or language constructors. Those decisions belong inside
module secrets and may change under the Module Guide's anticipated changes.

\section{Notation}

The structure of the MIS for modules comes from \citet{HoffmanAndStrooper1995},
with template modules adapted from \citet{GhezziEtAl2003}. The mathematical
notation follows Chapter~3 of \citet{HoffmanAndStrooper1995}. The symbol
$:=$ will be used for multiple assignment in Revision~0 semantics. Conditional
rules follow the form $(c_1 \Rightarrow r_1 \mid c_2 \Rightarrow r_2 \mid
\ldots \mid c_n \Rightarrow r_n)$.

\subsection{Primitive Data Types}

\begin{center}
\renewcommand{\arraystretch}{1.2}
\begin{tabular}{l l p{7.5cm}}
\toprule
\textbf{Data Type} & \textbf{Notation} & \textbf{Description}\\
\midrule
character & char & a single symbol or digit\\
integer & $\mathbb{Z}$ & a number without a fractional component\\
natural number & $\mathbb{N}$ & a number without a fractional component in
$[1, \infty)$\\
real & $\mathbb{R}$ & any number in $(-\infty, \infty)$\\
boolean & $\mathbb{B}$ & True or False\\
string & string & sequence of characters\\
\bottomrule
\end{tabular}
\end{center}

\subsection{Derived Types Used in Access Programs}

These types name conceptual values exchanged between modules. Their concrete
encoding is a secret of the owning module (typically Intent Model Store or Host
IDE Interface) and is not prescribed here.

\begin{center}
\renewcommand{\arraystretch}{1.2}
\begin{tabular}{l p{9.5cm}}
\toprule
\textbf{Type} & \textbf{Description}\\
\midrule
Request & A developer submission (natural language and/or supported artifact
reference).\\
RequestId & Opaque identifier of an in-progress workflow request.\\
Proposal & A draft set of requirements, constraints, and related items awaiting
decision.\\
Artifact & A stored intent-model entry (requirement, constraint, rationale, or
related record) as seen by clients.\\
ArtifactId & Identifier of an Artifact in the intent model.\\
ArtifactHistory & Ordered recoverable revisions of an Artifact.\\
Decision & One of $\{\mathrm{approve}, \mathrm{edit}, \mathrm{reject}\}$ with
optional edit payload.\\
Feedback & Developer refinement input applied to a Proposal.\\
Rationale & Developer-supplied justification recorded with an approval.\\
Note & Short developer-facing explanation (for example a revision request).\\
Stage & Workflow stage:
$\{\mathrm{capture}, \mathrm{review}, \mathrm{gate}, \mathrm{generate},
\mathrm{verify}\}$.\\
GateState & Whether generation is $\mathrm{blocked}$ or $\mathrm{released}$ for a
request.\\
PanelContent & Structured content presented on a host interactive surface
(encoding is a Host IDE Interface secret).\\
WorkspacePath & Reference to a workspace location relative to the open project
root (encoding is a Host IDE Interface secret).\\
WorkspaceContent & Contents obtained from a WorkspacePath.\\
Context & Prior approved or draft intent supplied to elicitation.\\
ContractSet & Compiled executable contracts derived from approved artifacts.\\
ImplCandidate & A candidate implementation (or change set) to be checked.\\
VerifyResult & Structured evidence of pass, fail, or inconclusive verification.\\
TraceLink & A justified association among artifacts, code paths, or tests.\\
LogEvent & A diagnostic record of a failure or significant interaction.\\
StatusMsg & A human-readable status or error message for the IDE surface.\\
SchemaId & Identifier of a supported intent-model schema.\\
\bottomrule
\end{tabular}
\end{center}

\section{Module Decomposition}

The following table is taken from the Module Guide for this project.

\begin{table}[H]
\centering
\begin{tabular}{p{0.32\textwidth} p{0.58\textwidth}}
\toprule
\textbf{Level 1} & \textbf{Level 2}\\
\midrule
Hardware-Hiding Module & Host IDE Interface Module\\
\midrule
Behaviour-Hiding Module & Workflow Controller Module\\
 & Intent Capture Module\\
 & Review Panel Module\\
 & Approval Gate Module\\
 & Verification Report Module\\
\midrule
Software Decision Module & Intent Model Store Module\\
 & Traceability Manager Module\\
 & Elicitation Engine Module\\
 & Contract Compiler Module\\
 & Contract Verifier Module\\
 & Coding Assistant Adapter Module\\
 & Diagnostic Log Module\\
\bottomrule
\end{tabular}
\caption{Module Hierarchy}
\label{TblMH}
\end{table}

Only leaf modules have MIS entries below. Non-leaf entries (Hardware-Hiding,
Behaviour-Hiding, Software Decision) are organizational and are not implemented.

%----------------------------------------------------------------------
\newpage
\section{MIS of Host IDE Interface Module} \label{MISHost}

\subsection{Module}

HostIDEInterface

\subsection{Uses}

None (uses host OS / IDE services as environment).

\subsection{Syntax}

\subsubsection{Exported Constants}

None.

\subsubsection{Exported Access Programs}

\begin{center}
\begin{tabular}{p{3.4cm} p{3.5cm} p{3.4cm} p{2.2cm}}
\hline
\textbf{Name} & \textbf{In} & \textbf{Out} & \textbf{Exceptions} \\
\hline
bindHostSession & - & - & HostUnavailable \\
presentContent & title: StatusMsg, body: PanelContent & - & HostUnavailable \\
presentStatus & msg: StatusMsg & - & - \\
promptDecision & proposal: Proposal & Decision & Cancelled, HostUnavailable \\
readWorkspace & path: WorkspacePath & WorkspaceContent & NotFound,
AccessDenied \\
observeWorkspace & - & - & HostUnavailable \\
\hline
\end{tabular}
\end{center}

\subsection{Semantics (Revision~-1 outline)}

\subsubsection{State Variables}

$bound$: $\mathbb{B}$ --- whether a host session is bound for this extension
activation.\\
$observed$: $\mathbb{B}$ --- whether workspace observation is active.

\subsubsection{Environment Variables}

$hostSurface$: the interactive command and presentation surface of the host IDE.\\
$workspaceRoot$: the open project folder in the host IDE.

\subsubsection{Assumptions}

Access programs other than \texttt{bindHostSession} are called only when
$bound = \mathrm{True}$.

\subsubsection{Access Routine Semantics}

Formal transition/output rules will be completed in Revision~0. For Rev~-1:
\texttt{bindHostSession} attaches to the host surface;
\texttt{presentContent} and \texttt{presentStatus} update the developer-visible
environment; \texttt{promptDecision} returns the developer's Decision or raises
\texttt{Cancelled}; workspace programs read or observe locations under
$workspaceRoot$. Host widget APIs remain secrets of this module.

\subsubsection{Local Functions}

None.

%----------------------------------------------------------------------
\newpage
\section{MIS of Workflow Controller Module} \label{MISControl}

\subsection{Module}

WorkflowController

\subsection{Uses}

IntentCapture, ReviewPanel, ApprovalGate, VerificationReport, HostIDEInterface,
DiagnosticLog

\subsection{Syntax}

\subsubsection{Exported Constants}

\texttt{STAGE\_ORDER}: seq of Stage --- the only legal stage succession for a
request.

\subsubsection{Exported Access Programs}

\begin{center}
\begin{tabular}{p{3.2cm} p{3.8cm} p{3.4cm} p{2.2cm}}
\hline
\textbf{Name} & \textbf{In} & \textbf{Out} & \textbf{Exceptions} \\
\hline
startRequest & req: Request & RequestId & InvalidRequest \\
advance & id: RequestId & Stage & IllegalStage, UnknownId \\
resumeDraft & id: RequestId & Stage & UnknownId, NotResumable \\
discardDraft & id: RequestId & - & UnknownId, ApprovedLocked \\
currentStage & id: RequestId & Stage & UnknownId \\
\hline
\end{tabular}
\end{center}

\subsection{Semantics (Revision~-1 outline)}

\subsubsection{State Variables}

$active$: map from RequestId to Stage.\\
$draftOpen$: map from RequestId to $\mathbb{B}$.

\subsubsection{Environment Variables}

None.

\subsubsection{Assumptions}

Callers invoke \texttt{startRequest} before other programs for a given request.
\texttt{discardDraft} does not remove already-approved intent held by the Intent
Model Store. \texttt{advance} may move only according to \texttt{STAGE\_ORDER}.

\subsubsection{Access Routine Semantics}

To be formalized in Revision~0. Stages follow capture $\rightarrow$ review
$\rightarrow$ gate $\rightarrow$ generate/verify as described in the Module
Guide.

\subsubsection{Local Functions}

None.

%----------------------------------------------------------------------
\newpage
\section{MIS of Intent Capture Module} \label{MISCapture}

\subsection{Module}

IntentCapture

\subsection{Uses}

ElicitationEngine, IntentModelStore, HostIDEInterface

\subsection{Syntax}

\subsubsection{Exported Constants}

None.

\subsubsection{Exported Access Programs}

\begin{center}
\begin{tabular}{p{3.2cm} p{3.8cm} p{3.4cm} p{2.2cm}}
\hline
\textbf{Name} & \textbf{In} & \textbf{Out} & \textbf{Exceptions} \\
\hline
submitRequest & req: Request & proposal: Proposal & UnsupportedArtifact,
InvalidRequest \\
refine & proposal: Proposal, feedback: Feedback & proposal: Proposal &
RefineFailed \\
saveDraft & proposal: Proposal & ArtifactId & StoreRejected \\
\hline
\end{tabular}
\end{center}

\subsection{Semantics (Revision~-1 outline)}

\subsubsection{State Variables}

$lastProposal$: Proposal $\cup \{\mathrm{nil}\}$.

\subsubsection{Environment Variables}

None.

\subsubsection{Assumptions}

Supported artifact kinds are those listed in the SRS intent-capture
requirements; unsupported kinds raise \texttt{UnsupportedArtifact}.

\subsubsection{Access Routine Semantics}

To be formalized in Revision~0. \texttt{submitRequest} and \texttt{refine}
delegate extraction/refinement to ElicitationEngine and return a draft Proposal
suitable for review. How natural language or artifacts become Proposal fields
is a secret of ElicitationEngine, not of this module's syntax.

\subsubsection{Local Functions}

None.

%----------------------------------------------------------------------
\newpage
\section{MIS of Review Panel Module} \label{MISReview}

\subsection{Module}

ReviewPanel

\subsection{Uses}

HostIDEInterface, IntentModelStore

\subsection{Syntax}

\subsubsection{Exported Constants}

None.

\subsubsection{Exported Access Programs}

\begin{center}
\begin{tabular}{p{3.2cm} p{3.8cm} p{3.4cm} p{2.2cm}}
\hline
\textbf{Name} & \textbf{In} & \textbf{Out} & \textbf{Exceptions} \\
\hline
present & proposal: Proposal & - & DisplayFailed \\
collectDecision & proposal: Proposal & Decision & Cancelled \\
amendItem & proposal: Proposal, item: ArtifactId, feedback: Feedback &
proposal: Proposal & UnknownItem \\
removeItem & proposal: Proposal, item: ArtifactId & proposal: Proposal &
UnknownItem \\
\hline
\end{tabular}
\end{center}

\subsection{Semantics (Revision~-1 outline)}

\subsubsection{State Variables}

$visibleProposal$: Proposal $\cup \{\mathrm{nil}\}$.

\subsubsection{Environment Variables}

None (presentation uses HostIDEInterface; the host surface is that module's
environment).

\subsubsection{Assumptions}

\texttt{collectDecision} is called only for a proposal that is currently
presented. Decisions originate from the developer via the host UI, not from the
coding assistant.

\subsubsection{Access Routine Semantics}

To be formalized in Revision~0.

\subsubsection{Local Functions}

None.

%----------------------------------------------------------------------
\newpage
\section{MIS of Approval Gate Module} \label{MISGate}

\subsection{Module}

ApprovalGate

\subsection{Uses}

IntentModelStore, CodingAssistantAdapter, ContractCompiler, HostIDEInterface,
DiagnosticLog

\subsection{Syntax}

\subsubsection{Exported Constants}

None.

\subsubsection{Exported Access Programs}

\begin{center}
\begin{tabular}{p{3.4cm} p{3.6cm} p{3.4cm} p{2.2cm}}
\hline
\textbf{Name} & \textbf{In} & \textbf{Out} & \textbf{Exceptions} \\
\hline
recordDecision & id: RequestId, d: Decision & GateState &
InvalidDecision, StoreFailed \\
generationAllowed & id: RequestId & $\mathbb{B}$ & UnknownId \\
releaseToAssistant & id: RequestId & - & Blocked, AssistFailed \\
requestRevision & id: RequestId, note: Note & - & AssistFailed \\
\hline
\end{tabular}
\end{center}

\subsection{Semantics (Revision~-1 outline)}

\subsubsection{State Variables}

$decision$: map from RequestId to Decision $\cup \{\mathrm{nil}\}$.\\
$gate$: map from RequestId to GateState.

\subsubsection{Environment Variables}

None.

\subsubsection{Assumptions}

\texttt{releaseToAssistant} requires a recorded $\mathrm{approve}$ decision.
On communication failure, the gate remains blocked and both the intent model and
assistant instruction are left unchanged (Hazard Analysis SR8).

\subsubsection{Access Routine Semantics}

To be formalized in Revision~0.

\subsubsection{Local Functions}

None.

%----------------------------------------------------------------------
\newpage
\section{MIS of Verification Report Module} \label{MISReport}

\subsection{Module}

VerificationReport

\subsection{Uses}

ContractVerifier, HostIDEInterface

\subsection{Syntax}

\subsubsection{Exported Constants}

None.

\subsubsection{Exported Access Programs}

\begin{center}
\begin{tabular}{p{3.2cm} p{3.8cm} p{3.4cm} p{2.2cm}}
\hline
\textbf{Name} & \textbf{In} & \textbf{Out} & \textbf{Exceptions} \\
\hline
runAndPresent & candidate: ImplCandidate, id: RequestId &
VerifyResult & VerifyFailed \\
present & result: VerifyResult & - & DisplayFailed \\
\hline
\end{tabular}
\end{center}

\subsection{Semantics (Revision~-1 outline)}

\subsubsection{State Variables}

$lastResult$: VerifyResult $\cup \{\mathrm{nil}\}$.

\subsubsection{Environment Variables}

None.

\subsubsection{Assumptions}

An inconclusive result from the verifier is presented as inconclusive, never as
pass (Hazard Analysis SR11).

\subsubsection{Access Routine Semantics}

To be formalized in Revision~0.

\subsubsection{Local Functions}

None.

%----------------------------------------------------------------------
\newpage
\section{MIS of Intent Model Store Module} \label{MISStore}

\subsection{Module}

IntentModelStore

\subsection{Uses}

None.

\subsection{Syntax}

\subsubsection{Exported Constants}

\texttt{SUPPORTED\_SCHEMA}: SchemaId --- identifier of the schema this store
accepts for writes.

\subsubsection{Exported Access Programs}

\begin{center}
\begin{tabular}{p{3.4cm} p{3.6cm} p{3.4cm} p{2.2cm}}
\hline
\textbf{Name} & \textbf{In} & \textbf{Out} & \textbf{Exceptions} \\
\hline
openWorkspaceModel & - & - & AccessDenied \\
putDraft & proposal: Proposal & ArtifactId & SchemaInvalid \\
putApproved & proposal: Proposal, rationale: Rationale & seq of ArtifactId &
SchemaInvalid \\
get & id: ArtifactId & Artifact & NotFound \\
history & id: ArtifactId & ArtifactHistory & NotFound \\
validate & proposal: Proposal & $\mathbb{B}$ & - \\
\hline
\end{tabular}
\end{center}

\subsection{Semantics (Revision~-1 outline)}

\subsubsection{State Variables}

$model$: the set of stored artifacts and drafts.\\
$open$: $\mathbb{B}$ --- whether the workspace model is open for this session.

\subsubsection{Environment Variables}

$workspaceStore$: the workspace-backed store that holds persisted intent-model
entries (Git-friendly project storage). Path layout and encoding are secrets.

\subsubsection{Assumptions}

Writers call \texttt{openWorkspaceModel} once per workspace session before other
programs. Schema-invalid writes raise \texttt{SchemaInvalid} and leave $model$
unchanged. Clients receive Artifact values, never raw encodings.

\subsubsection{Access Routine Semantics}

To be formalized in Revision~0. Persistence layout is a secret of this module
(Module Guide anticipated change for schema).

\subsubsection{Local Functions}

None.

%----------------------------------------------------------------------
\newpage
\section{MIS of Traceability Manager Module} \label{MISTrace}

\subsection{Module}

TraceabilityManager

\subsection{Uses}

IntentModelStore

\subsection{Syntax}

\subsubsection{Exported Constants}

None.

\subsubsection{Exported Access Programs}

\begin{center}
\begin{tabular}{p{3.2cm} p{3.8cm} p{3.4cm} p{2.2cm}}
\hline
\textbf{Name} & \textbf{In} & \textbf{Out} & \textbf{Exceptions} \\
\hline
addLink & link: TraceLink & - & InvalidLink, NotFound \\
linksFrom & id: ArtifactId & seq of TraceLink & NotFound \\
linksTo & id: ArtifactId & seq of TraceLink & NotFound \\
reportBroken & - & seq of TraceLink & - \\
\hline
\end{tabular}
\end{center}

\subsection{Semantics (Revision~-1 outline)}

\subsubsection{State Variables}

$links$: set of TraceLink.

\subsubsection{Environment Variables}

None.

\subsubsection{Assumptions}

Endpoints of a link must exist in the Intent Model Store (or denote an allowed
code/test reference form defined at Revision~0). How link justifications are
stored is a secret of this module.

\subsubsection{Access Routine Semantics}

To be formalized in Revision~0. \texttt{reportBroken} returns missing, broken,
or unjustified links (Hazard Analysis SR10).

\subsubsection{Local Functions}

None.

%----------------------------------------------------------------------
\newpage
\section{MIS of Elicitation Engine Module} \label{MISElicit}

\subsection{Module}

ElicitationEngine

\subsection{Uses}

IntentModelStore, DiagnosticLog

\subsection{Syntax}

\subsubsection{Exported Constants}

\texttt{MAX\_CONTEXT\_UNITS}: $\mathbb{N}$ --- upper bound on Context size
accepted by \texttt{extract} before raising \texttt{BoundExceeded}.

\subsubsection{Exported Access Programs}

\begin{center}
\begin{tabular}{p{3.2cm} p{3.8cm} p{3.4cm} p{2.2cm}}
\hline
\textbf{Name} & \textbf{In} & \textbf{Out} & \textbf{Exceptions} \\
\hline
extract & req: Request, context: Context & proposal: Proposal &
ModelUnavailable, BoundExceeded \\
applyFeedback & proposal: Proposal, feedback: Feedback & proposal: Proposal &
ModelUnavailable \\
\hline
\end{tabular}
\end{center}

\subsection{Semantics (Revision~-1 outline)}

\subsubsection{State Variables}

None required beyond transient computation.

\subsubsection{Environment Variables}

$llmService$: external language-model service used for extraction (optional at
runtime; unavailability raises \texttt{ModelUnavailable}). Prompt templates and
orchestration structure are secrets of this module, not exported constants.

\subsubsection{Assumptions}

Egress of project content to $llmService$ respects the privacy policy of Hazard
Analysis SR12. Oversized inputs raise \texttt{BoundExceeded}.

\subsubsection{Access Routine Semantics}

To be formalized in Revision~0.

\subsubsection{Local Functions}

None.

%----------------------------------------------------------------------
\newpage
\section{MIS of Contract Compiler Module} \label{MISCompile}

\subsection{Module}

ContractCompiler

\subsection{Uses}

IntentModelStore, TraceabilityManager

\subsection{Syntax}

\subsubsection{Exported Constants}

None.

\subsubsection{Exported Access Programs}

\begin{center}
\begin{tabular}{p{3.2cm} p{3.8cm} p{3.4cm} p{2.2cm}}
\hline
\textbf{Name} & \textbf{In} & \textbf{Out} & \textbf{Exceptions} \\
\hline
compile & ids: seq of ArtifactId & ContractSet & IncompleteIntent,
CompileFailed \\
\hline
\end{tabular}
\end{center}

\subsection{Semantics (Revision~-1 outline)}

\subsubsection{State Variables}

None.

\subsubsection{Environment Variables}

None.

\subsubsection{Assumptions}

All ArtifactIds passed to \texttt{compile} refer to approved artifacts. The
compiled ContractSet remains associated with those source ids through the
Traceability Manager.

\subsubsection{Access Routine Semantics}

To be formalized in Revision~0. Choice of contract notation and the mapping from
intent artifacts to checkable obligations are secrets of this module (Module
Guide anticipated change for contracts).

\subsubsection{Local Functions}

None.

%----------------------------------------------------------------------
\newpage
\section{MIS of Contract Verifier Module} \label{MISVerify}

\subsection{Module}

ContractVerifier

\subsection{Uses}

IntentModelStore, TraceabilityManager, DiagnosticLog

\subsection{Syntax}

\subsubsection{Exported Constants}

None.

\subsubsection{Exported Access Programs}

\begin{center}
\begin{tabular}{p{3.2cm} p{3.8cm} p{3.4cm} p{2.2cm}}
\hline
\textbf{Name} & \textbf{In} & \textbf{Out} & \textbf{Exceptions} \\
\hline
check & candidate: ImplCandidate, contracts: ContractSet & VerifyResult &
CheckerUnavailable \\
\hline
\end{tabular}
\end{center}

\subsection{Semantics (Revision~-1 outline)}

\subsubsection{State Variables}

None.

\subsubsection{Environment Variables}

$checker$: external contract-checking tool available to this module. Selection
and invocation details are secrets.

\subsubsection{Assumptions}

If $checker$ is unavailable or fails infrastructurally, the result is
inconclusive and \texttt{CheckerUnavailable} may be raised; this must not be
interpreted as pass.

\subsubsection{Access Routine Semantics}

To be formalized in Revision~0.

\subsubsection{Local Functions}

None.

%----------------------------------------------------------------------
\newpage
\section{MIS of Coding Assistant Adapter Module} \label{MISAssist}

\subsection{Module}

CodingAssistantAdapter

\subsection{Uses}

IntentModelStore, DiagnosticLog

\subsection{Syntax}

\subsubsection{Exported Constants}

None.

\subsubsection{Exported Access Programs}

\begin{center}
\begin{tabular}{p{3.6cm} p{3.4cm} p{3.4cm} p{2.2cm}}
\hline
\textbf{Name} & \textbf{In} & \textbf{Out} & \textbf{Exceptions} \\
\hline
sendApproved & id: RequestId, ids: seq of ArtifactId & - &
NotApproved, ChannelFailed \\
sendRevisionRequest & id: RequestId, note: Note & - & ChannelFailed \\
receiveProposal & - & ImplCandidate & Timeout, ChannelFailed \\
\hline
\end{tabular}
\end{center}

\subsection{Semantics (Revision~-1 outline)}

\subsubsection{State Variables}

$connected$: $\mathbb{B}$.

\subsubsection{Environment Variables}

$assistantChannel$: communication channel to the AI coding assistant. Message
encoding and tool registration are secrets of this module.

\subsubsection{Assumptions}

\texttt{sendApproved} is invoked only when the Approval Gate has released
generation. Approve/reject are not exposed as assistant-callable operations
(Hazard Analysis SR7).

\subsubsection{Access Routine Semantics}

To be formalized in Revision~0.

\subsubsection{Local Functions}

None.

%----------------------------------------------------------------------
\newpage
\section{MIS of Diagnostic Log Module} \label{MISLog}

\subsection{Module}

DiagnosticLog

\subsection{Uses}

None.

\subsection{Syntax}

\subsubsection{Exported Constants}

\texttt{MAX\_RECENT}: $\mathbb{N}$ --- maximum length returned by
\texttt{recent}.

\subsubsection{Exported Access Programs}

\begin{center}
\begin{tabular}{p{3.2cm} p{3.8cm} p{3.4cm} p{2.2cm}}
\hline
\textbf{Name} & \textbf{In} & \textbf{Out} & \textbf{Exceptions} \\
\hline
record & event: LogEvent & - & - \\
recent & n: $\mathbb{N}$ & seq of LogEvent & - \\
\hline
\end{tabular}
\end{center}

\subsection{Semantics (Revision~-1 outline)}

\subsubsection{State Variables}

$events$: seq of LogEvent.

\subsubsection{Environment Variables}

$logSink$: optional persistent sink for diagnostic records; location and format
are secrets of this module.

\subsubsection{Assumptions}

Logging must not record secret material from Git-ignored content (aligned with
HA-SR12). If $n > \texttt{MAX\_RECENT}$, \texttt{recent} returns at most
\texttt{MAX\_RECENT} events.

\subsubsection{Access Routine Semantics}

To be formalized in Revision~0.

\subsubsection{Local Functions}

None.

%----------------------------------------------------------------------
\newpage

\section{Appendix --- Generative AI Use Disclosure}

\subsection{AI Tools Used}

\begin{itemize}
  \item Cursor and Claude, used in the writing process for this MIS.
\end{itemize}

\subsection{How and Where Used}

\begin{itemize}
  \item \textbf{Completing sections from drafts.} Used to complete and fill
    out MIS sections based on the author's own drafts of module interfaces,
    types, and access programs.
  \item \textbf{Diagrams and tables.} Used to help format and produce syntax
    tables and related structured content for the module specifications.
  \item \textbf{Rubric review.} Used to evaluate the document against the
    Design Documentation rubric and to suggest areas of improvement, which
    the author then accepted, revised, or rejected.
  \item \textbf{Reflection.} Vedant Patel's reflection answers were drafted
    with Claude from his own account of the deliverable, then reviewed and
    edited by him.
\end{itemize}

\subsection{Verification of AI-Generated Content}

\begin{enumerate}
  \item Reviewed all AI-assisted expansions against the author's drafts and
    the Module Guide hierarchy before keeping them.
  \item Checked that exported constants, access programs, state variables, and
    environment variables are present for each leaf module (with explicit
    ``None'' where empty).
  \item Spot-checked that Uses lists remain consistent with the Module Guide
    and that Formal $:=$ access-routine semantics remain deferred to
    Revision~0 per Des.\ Doc.\ Rev~-1 scope (syntax only).
\end{enumerate}

These steps are intended to convince a reader of the accuracy, completeness,
consistency, and appropriateness of this MIS. The team remains responsible for
the document.

\newpage

\section{Appendix --- Reflection}

The information in this section will be used to evaluate the team members on the
graduate attribute of Problem Analysis and Design.

\input{../../Reflection.text}

\begin{enumerate}
  \item What went well while writing this deliverable?

  \textbf{Vedant Patel:}

  Having the Module Guide and MIS drafted early, on November 13, let the rest of
  us work on the proof of concept and the Wireframe Report in parallel instead
  of waiting on the design. Building the wireframe was also a useful check on
  the design. Every screen had to be backed by a module, so I could confirm that
  the Review Panel, Approval Gate, Traceability Manager, and Verification Report
  modules covered everything a developer actually sees. The proof-of-concept
  code helped in the same way: the intent capture and stakeholder states I
  worked with in the Goal2Story stage matched the data the Elicitation Engine
  and Intent Model Store are specified to hold.

  \item What pain points did you experience during this deliverable, and how
    did you resolve them?

  \textbf{Vedant Patel:}

  The biggest pain point was confusion over which custom document we owed. Our
  Problem Statement and VnV Plan said the fall CD was a Literature Review, but
  the course roster, which is the official record, listed a Wireframe Report.
  We only noticed the mismatch close to the deadline. I resolved it by checking
  the course repository directly, confirming the registered CDs and the due
  date, and updating our documents to match. A second pain point was our build
  pipeline. Merges to main kept failing, and we found that the CI could not push
  the generated PDFs because of a branch protection rule, so the published PDFs
  were out of date. Knowing the cause meant we could plan a fix instead of
  guessing. The lesson for me is to check the official source for requirements
  and deadlines early, not our own copies of them.

  \item Which of your design decisions stemmed from speaking to your client(s)
  or a proxy (e.g. your peers, stakeholders, potential users)? For those that
  were not, why, and where did they come from?

  \textbf{Vedant Patel:}

  The clearest input from proxies this round came from showing the wireframe to
  classmates. They found the screens text heavy and wanted an introduction
  explaining what the tool is for and how to use it. Both points affect the
  Review Panel and Host IDE Interface modules: the Review Panel should lead with
  the decision the developer has to make, and the Host IDE Interface needs a
  first-run walkthrough. Most other design decisions came from the SRS, which
  was itself based on the project description from our supervisor. For
  example, the Approval Gate exists because the SRS requires that no code is
  generated before the developer approves the requirements (FR7). We relied on
  the SRS because it already captured what our stakeholders asked for, and
  there was not time to validate every module directly with users before Rev~-1.

  \item While creating the design doc, what parts of your other documents (e.g.
  requirements, hazard analysis, etc), it any, needed to be changed, and why?

  \textbf{Vedant Patel:}

  The custom document sections of the Problem Statement and the VnV Plan had to
  change from Literature Review and Code Walkthroughs to Wireframe Report and
  Usability Report, so they matched what the team registered with the course.
  I also corrected the proof-of-concept dates in the SRS schedule constraints,
  which listed November 30 to December 3 instead of the course window of
  November 23 to December 4. Neither change affected the requirements
  themselves, but both would have misled a reader about our plan.

  \item What are the limitations of your solution?  Put another way, given
  unlimited resources, what could you do to make the project better? (LO\_ProbSolutions)

  \textbf{Vedant Patel:}

  Our design depends heavily on the quality of the language model output. If
  the Elicitation Engine proposes poor stakeholders or user stories, the
  developer has to catch it during review, and the classmate feedback suggests
  that review already feels text heavy. Traceability is also only as good as the
  coding assistant's report of which files it changed, and verification only
  covers what the compiled contracts can express, which for now means Java with
  JML. With unlimited resources I would run proper user studies with developers
  who use AI assistants daily, support more languages and IDEs, add automatic
  link recovery as a backup when the assistant's report is incomplete, and
  measure how much time the review steps really cost compared with the time
  they save later.

  \item Give a brief overview of other design solutions you considered.  What
  are the benefits and tradeoffs of those other designs compared with the chosen
  design?  From all the potential options, why did you select the documented design?
  (LO\_Explores)

  \textbf{Vedant Patel:}

  One alternative was to recover traceability after the fact, by analysing code
  and documents to guess which requirement each piece of code implements. That
  needs nothing from the developer up front, but the links are only guesses and
  the reasons behind a decision cannot be recovered at all. We chose to record
  links and rationale at the moment the developer approves a requirement,
  because that is when the information exists. For the interface, I compared a
  standalone web application with panels inside the IDE. A web app would be
  easier to build, but the SRS requires the tool to work inside the IDE without
  a separate interface (ACCR1), and developers are less likely to switch
  windows mid-task. For intent capture, I compared a structured form with a
  conversation. A form is faster for experienced users, but a conversation with
  suggested answers is easier to start with, so the wireframe uses that. In the
  proof of concept we also split elicitation into separate stages (capture
  intent, stakeholders, goals to user stories) instead of one large prompt,
  because each stage can be checked and approved on its own, which matches the
  modular structure in the Module Guide.

  \item (After you have implemented another team's module, which means this
  isn't filled in until after the original deadline). What did you learn by
  implementing another team's module? Were all the details you needed in the
  documentation, or did you need to make assumptions, or ask the other team
  questions?  If your team also had another team implement one of your modules,
  what was this experience like?  Are there things in your documentation you
  could have changed to make the process go more smoothly for when an
  ``outsider'' completes some of the implementation?

  To be completed after implementing another team's module, as scheduled for
  Design Document Revision~0.

\end{enumerate}

\bibliographystyle{plainnat}
\bibliography{../../../refs/References}

\end{document}
